\section{Билет 65. Второе начало термодинамики в формулировке Кельвина и в формулировке Клаузиуса. Необратимый цикл. Неравенство Клаузиуса.}

\begin{definition}
\textbf{Второе начало термодинамики (формулировка Кельвина):} Невозможно построить периодически действующую тепловую машину, единственным результатом работы которой было бы превращение всего количества теплоты, полученного от нагревателя, в механическую работу. Иными словами, вечный двигатель второго рода невозможен. Любая тепловая машина должна иметь не только нагреватель, но и холодильник, которому часть энергии неизбежно передается в виде теплоты.
\end{definition}

\begin{definition}
\textbf{Второе начало термодинамики (формулировка Клаузиуса):} Невозможен процесс, единственным результатом которого была бы передача теплоты от менее нагретого тела к более нагретому без совершения работы или других изменений в системе. Самопроизвольный переход тепла от холодного тела к горячему запрещен; для этого требуется внешний источник энергии (как в холодильной машине).
\end{definition}

\begin{theorem}[Неравенство Клаузиуса]
Для любого кругового процесса (обратимого или необратимого), совершаемого термодинамической системой, интеграл от приведенной теплоты удовлетворяет соотношению:
\begin{equation}
    \oint \frac{\delta Q}{T} \le 0, \label{eq:clausius_ineq_65}
\end{equation}
где $\delta Q$ --- элементарное количество теплоты, полученное системой при температуре $T$.
\begin{itemize}
    \item Знак ``$=$'' относится к \textbf{обратимым} циклам (например, циклу Карно).
    \item Знак ``$<$'' относится к \textbf{необратимым} (реальным) циклам.
\end{itemize}
\end{theorem}

\begin{property}[Вывод неравенства для необратимого цикла]
Рассмотрим реальную тепловую машину, работающую по необратимому циклу между температурами $T_1$ и $T_2$. Согласно теореме Карно, её КПД строго меньше КПД идеальной машины Карно: $\eta_{\text{необр}} < \eta_{\text{К}}$.
\begin{equation}
    1 - \frac{|Q_2|}{Q_1} < 1 - \frac{T_2}{T_1} \quad \Rightarrow \quad \frac{Q_1}{T_1} < \frac{|Q_2|}{T_2}.
\end{equation}
Учитывая знаки теплот ($Q_1 > 0$, $Q_2 < 0$), получаем:
\begin{equation}
    \frac{Q_1}{T_1} + \frac{Q_2}{T_2} < 0 \quad \Rightarrow \quad \oint \frac{\delta Q}{T} < 0.
\end{equation}
Это неравенство обобщается на любые необратимые процессы и выражает фундаментальную направленность реальных термодинамических процессов.
\end{property}

\begin{remark}
\textbf{Физический смысл:} Неравенство Клаузиуса математически выражает необратимость реальных процессов (трение, теплопроводность, диффузия, неупругие деформации). В изолированной системе ($\delta Q = 0$) из него следует $\Delta S > 0$, что означает монотонный рост энтропии до достижения термодинамического равновесия. Равновесному состоянию соответствует максимум энтропии.
\end{remark}
